\chapter{Πλήρεις ανακατασκευές σημάτων Ringdown}
\label{app:ringdown-reconstructions}

Το παρόν παράρτημα περιλαμβάνει το πλήρες σύνολο των διαγραμμάτων
ανακατασκευής που προέκυψαν από την εφαρμογή της μεθόδου Matrix Pencil.
Για κάθε γεννήτρια παρουσιάζονται τα τέσσερα διαθέσιμα σήματα: τάση
($V$), ρεύμα ($I$), ενεργός ισχύς ($P$) και άεργος ισχύς ($Q$). Κάθε
διάγραμμα περιλαμβάνει τις εξεταζόμενες προσεγγίσεις και επιτρέπει την
άμεση σύγκριση του αρχικού φιλτραρισμένου σήματος με την αντίστοιχη
ανακατασκευή.

\newcommand{\ringdownreconstructionfigure}[5]{%
\begin{figure}[p]
    \centering
    \includegraphics[width=0.94\textwidth,height=0.80\textheight,keepaspectratio]{#1#2}
    \caption{Ανακατασκευή σήματος #3 για τη γεννήτρια $G_{#4}$ του #5.}
\end{figure}%
}

\section{Kundur Two-Area System}

Τα ακόλουθα σχήματα παρουσιάζουν τις ανακατασκευές για όλες τις
γεννήτριες του συστήματος Kundur Two-Area.

\foreach \gen in {1,2,3,4}{%
    \ringdownreconstructionfigure{PreliminaryInvestigation/plots/reconstruction_grids/pdf/}{g\gen\string_Active\string_Power\string_Reconstruction.pdf}{ενεργού ισχύος ($P$)}{\gen}{\textit{Kundur Two-Area System}}%
    \ringdownreconstructionfigure{PreliminaryInvestigation/plots/reconstruction_grids/pdf/}{g\gen\string_Current\string_Reconstruction.pdf}{ρεύματος ($I$)}{\gen}{\textit{Kundur Two-Area System}}%
    \ringdownreconstructionfigure{PreliminaryInvestigation/plots/reconstruction_grids/pdf/}{g\gen\string_Reactive\string_Power\string_Reconstruction.pdf}{αέργου ισχύος ($Q$)}{\gen}{\textit{Kundur Two-Area System}}%
    \ringdownreconstructionfigure{PreliminaryInvestigation/plots/reconstruction_grids/pdf/}{g\gen\string_Voltage\string_Reconstruction.pdf}{τάσης ($V$)}{\gen}{\textit{Kundur Two-Area System}}%
}

\clearpage
\section{IEEE 39-Bus System}

Τα ακόλουθα σχήματα παρουσιάζουν τις ανακατασκευές για όλες τις
γεννήτριες του συστήματος IEEE 39-bus.

\foreach \gen in {1,2,3,4,5,6,7,8,9,10}{%
    \ringdownreconstructionfigure{IEEE39/analysis/Load03_Pplus2_50s_0.2_to_end_reset/plots/reconstruction_grids/pdf/}{g\gen\string_Active\string_Power\string_reconstruction.pdf}{ενεργού ισχύος ($P$)}{\gen}{\textit{IEEE 39-Bus System}}%
    \ringdownreconstructionfigure{IEEE39/analysis/Load03_Pplus2_50s_0.2_to_end_reset/plots/reconstruction_grids/pdf/}{g\gen\string_Current\string_reconstruction.pdf}{ρεύματος ($I$)}{\gen}{\textit{IEEE 39-Bus System}}%
    \ringdownreconstructionfigure{IEEE39/analysis/Load03_Pplus2_50s_0.2_to_end_reset/plots/reconstruction_grids/pdf/}{g\gen\string_Reactive\string_Power\string_reconstruction.pdf}{αέργου ισχύος ($Q$)}{\gen}{\textit{IEEE 39-Bus System}}%
    \ringdownreconstructionfigure{IEEE39/analysis/Load03_Pplus2_50s_0.2_to_end_reset/plots/reconstruction_grids/pdf/}{g\gen\string_Voltage\string_reconstruction.pdf}{τάσης ($V$)}{\gen}{\textit{IEEE 39-Bus System}}%
}
